\documentclass[9pt,a4paper]{article}

\usepackage[margin=1.2cm]{geometry}
\usepackage[table]{xcolor}
\usepackage{array}
\usepackage{graphicx}
\usepackage{booktabs}
\usepackage[T1]{fontenc}
\usepackage[utf8]{inputenc}
\usepackage[hidelinks]{hyperref}
\usepackage{caption}
\usepackage{ragged2e}

\pagestyle{empty}

\newcolumntype{L}[1]{>{\RaggedRight\arraybackslash}p{#1}}

\begin{document}

\begin{table}[t]
\centering
\scriptsize
\setlength{\tabcolsep}{4pt}
\renewcommand{\arraystretch}{1.1}
\resizebox{\linewidth}{!}{%
\begin{tabular}{|L{0.47\linewidth}|L{0.47\linewidth}|}
\hline
\multicolumn{2}{|c|}{%
  \cellcolor[HTML]{A0A0A0}%
  \textbf{NORMA-O Requirements Specification
  Document}} \\
\hline

%------------------------------------------------------------
% 1. Purpose
%------------------------------------------------------------
\multicolumn{2}{|c|}{%
  \cellcolor[HTML]{EFEFEF}\textbf{1. Purpose}} \\
\hline
\multicolumn{2}{|L{0.96\linewidth}|}{%
NORMA-O is an OWL~2~DL ontology for representing legal provisions as formal norm instances. Each instance captures both the normative content (deontic type, legal agent, regulated action, legal object, and source provision) and the extraction context, including who produced the annotation, by what method, with what confidence, and under what review status.} \\
\hline

%------------------------------------------------------------
% 2. Scope
%------------------------------------------------------------
\multicolumn{2}{|c|}{%
  \cellcolor[HTML]{EFEFEF}\textbf{2. Scope}} \\
\hline
\multicolumn{2}{|L{0.96\linewidth}|}{%
NORMA-O covers legal norms extracted from regulatory provisions, including their deontic type, structural arguments, and conditional applicability. It also represents the annotation context that produced each norm instance: annotator, extraction method, confidence, and review status} \\
\hline

%------------------------------------------------------------
% 3. Implementation Language
%------------------------------------------------------------
\multicolumn{2}{|c|}{%
  \cellcolor[HTML]{EFEFEF}%
  \textbf{3. Implementation Language}} \\
\hline
\multicolumn{2}{|L{0.96\linewidth}|}{%
OWL~2, serialised in Turtle.} \\
\hline

%------------------------------------------------------------
% 4. Intended End-Users
%------------------------------------------------------------
\multicolumn{2}{|c|}{%
  \cellcolor[HTML]{EFEFEF}\textbf{4. Intended End-Users}} \\
\hline
\multicolumn{2}{|L{0.96\linewidth}|}{%
\textbf{U1.} Domain experts involved in regulatory compliance
work.} \\
\hline

%------------------------------------------------------------
% 5. Intended Uses
%------------------------------------------------------------
\multicolumn{2}{|c|}{%
  \cellcolor[HTML]{EFEFEF}\textbf{5. Intended Uses}} \\
\hline
\multicolumn{2}{|L{0.96\linewidth}|}{%
\textbf{Use~1.} Formal schema for representing and querying
legal norm instances.
\textbf{Use~2.} Vocabulary for deriving applicable norms via
SWRL rules under factual conditions.
\textbf{Use~3.} Foundation for SPARQL queries
over regulatory annotations.
\textbf{Use~4.} Reusable model applicable to any regulatory
domain.} \\
\hline

%------------------------------------------------------------
% 6a. Non-Functional Requirements
%------------------------------------------------------------
\multicolumn{2}{|c|}{%
  \cellcolor[HTML]{EFEFEF}%
  \textbf{6. Ontology Requirements}} \\
\hline
\multicolumn{2}{|c|}{%
  \cellcolor[HTML]{EFEFEF}%
  \textbf{6a. Non-Functional Requirements}} \\
\hline
\multicolumn{2}{|L{0.96\linewidth}|}{%
\textbf{NFR1.} FAIR compliant
\textbf{NFR2.} Consistent under OWL~2; no unsatisfiable classes.
\textbf{NFR3.} Extensible without modifying existing definitions.} \\
\hline

%------------------------------------------------------------
% 6b. Competency Questions
%------------------------------------------------------------
\multicolumn{2}{|c|}{%
  \cellcolor[HTML]{EFEFEF}%
  \textbf{6b. Functional Requirements --- Competency Questions
  (aligned with Section~7 of the paper)}} \\
\hline

\cellcolor[HTML]{EFEFEF}%
\textbf{CQG1. Normative content retrieval} &
\cellcolor[HTML]{EFEFEF}%
\textbf{CQG2. Source traceability and provenance} \\
\hline

CQ1. Which norms apply to a given legal agent?

CQ2. Which norms concern a given legal object?

CQ3. Which norms regulate a given legal action?

CQ4. Which norms originate from a given regulation, article, or
paragraph?

CQ5. Which source text is linked to a given norm?

CQ6. Which norms share the same legal source? &

CQ7. What annotator, extraction method, and confidence score are
recorded for each norm?

CQ8. Which norms carry a given review status?

CQ9. Which norms carry a given binding force, norm status, or
compliance criticality?

CQ10. Which norms are classified as critical or high compliance
priority?

CQ11. Which norms are derivable given a factual scenario?
(Answered via SWRL inference; requires a compatible reasoner.)\\
\hline

\end{tabular}%
}
\caption{NORMA-O Requirements Specification Document,
  following the LOT methodology template.
  \url{https://w3id.org/def/norma-o}.}
\label{tab:orsd}
\end{table}

\end{document}